% chapters/15/section-06.tex
\section{Dari Pygame ke Unity, dan Hak Pakai Bunyi}
\label{sec:dari-pygame}

\subsection{Bandingkan dengan Cara yang Sudah Anda Pakai}

Pada Bab~III sampai Bab~XI Anda memakai Pygame, dan bunyi di sana diurus dua
modul: \code{pygame.mixer.Sound} untuk efek suara dan \code{pygame.mixer.music}
untuk musik latar \cite{sweigart-invent}. Pemisahan itu bukan kebetulan, melainkan
kebutuhan yang sama yang Anda temui kembali pada Unity. Perbandingannya
diperlihatkan pada Tabel~\ref{tab:banding-bunyi}. Awalan \code{pygame.} pada
kolom kedua ditinggalkan agar tabelnya muat.

\begin{table}[htbp]
\centering
\caption{Pengurusan bunyi pada Pygame dan pada Unity}
\label{tab:banding-bunyi}
{\setlength{\tabcolsep}{4pt}\footnotesize%
\begin{tabularx}{\textwidth}{@{}L{3.5cm} L{4.4cm} Y@{}}
\toprule
\headrow \textbf{Keperluan} & \textbf{Pygame} & \textbf{Unity} \\
\midrule
Memuat berkas bunyi & \code{mixer.Sound("a.wav")} &
Diimpor seperti gambar, lalu diseret ke peubah \code{AudioClip} pada
\istilah{Inspector}. \\
\addlinespace
Memainkan efek suara & \code{bunyi.play()} &
\code{sfxSource.PlayOneShot(klip)}. \\
\addlinespace
Memainkan musik berulang & \code{mixer.music.load()} lalu
\code{mixer.music.play(-1)} &
\code{musicSource.loop = true} lalu \code{musicSource.Play()}. \\
\addlinespace
Menghentikan musik & \code{mixer.music.stop()} &
\code{musicSource.Stop()}. \\
\addlinespace
Mengatur kekuatan & \code{bunyi.set_volume(0.5)} &
Nilai pada \istilah{AudioSource}, atau tingkat pada \istilah{AudioMixer}. \\
\addlinespace
Melemahkan menurut jarak & Harus dihitung sendiri dari jarak pemain &
Diurus mesin melalui \istilah{Spatial Blend} dan \istilah{Min Distance}. \\
\bottomrule
\end{tabularx}}
\end{table}

Perhatikan baris terakhir. Di Pygame, melemahkan bunyi menurut jarak menuntut
Anda menghitung jaraknya sendiri, lalu memanggil \code{set_volume} setiap
bingkai. Di Unity, pekerjaan itu diurus mesin, dan Anda hanya perlu menggeser
satu nilai. Itulah pola yang berulang sepanjang Bab~XII sampai bab ini: Unity
memindahkan pekerjaan yang berulang ke pengaturan, sehingga Anda dapat
memusatkan perhatian pada hal yang khas bagi permainan Anda.

Satu pelajaran dari Pygame tetap berlaku, dan sayangnya sering dilupakan.
Sediakan jalan bagi pemain untuk mematikan bunyinya. Pada Pygame, caranya
memelihara satu peubah yang dibalik setiap kali tombol ditekan. Pada Unity,
caranya serupa: satu peubah yang memutuskan apakah musik dan efek suara
dipanggil \cite{stmikjayakarta-pygame}. Pemain yang bermain di tempat ramai,
atau yang sedang tidak ingin terdengar, akan berterima kasih untuk pilihan itu.

\subsection{Hak Pakai Berkas Bunyi}

Aturan yang dibahas pada Bab~XIV berlaku sama untuk bunyi, dan akibatnya justru
lebih berat. Berkas musik lebih mudah dikenali daripada gambar, penuntutan hak
ciptanya lebih lazim, dan satu potongan musik yang dipakai tanpa izin dapat
membatalkan seluruh proyek yang sedang Anda kerjakan.

\begin{itemize}
  \item \textbf{Cari tempat yang menyatakan lisensinya dengan jelas.} Situs
        yang menyediakan musik bebas biasanya mencantumkan jenis lisensinya
        pada setiap berkas.
  \item \textbf{Catat sejak mengunduh}, bukan sejak lupa. Catat nama berkas,
        sumbernya, jenis lisensinya, dan apakah penciptanya menuntut namanya
        dicantumkan.
  \item \textbf{Pisahkan musik untuk keperluan belajar dari musik untuk
        keperluan yang lebih luas.} Sebagian lisensi mengizinkan pemakaian
        untuk belajar, tetapi tidak mengizinkan pemakaian dalam karya yang
        dijual. Periksa hal itu sebelum Anda menyebarkan permainannya.
\end{itemize}

\begin{tipsbox}
Bunyi juga dapat dibuat sendiri, dan untuk keperluan mata kuliah ini cara itu
sering lebih cepat daripada mencari. Sebuah bunyi lompat dapat dibuat dari satu
bunyi pendek yang nadanya dinaikkan sedikit. Buku ini sendiri memakai kebiasaan
mencatat sumber: setiap gambar yang bukan karya penulis diberi keterangan
sumbernya di bawah gambarnya, misalnya pada
Gambar~\ref{fig:petunjuk-visual-dan-bunyi}.
\end{tipsbox}

\subsection{Ringkasan Perkembangan}

\begin{table}[htbp]
\centering
\caption{Bagian yang ditambahkan pada bunyi dan umpan balik pada bab berikutnya}
\label{tab:perkembangan-bunyi}
{\setlength{\tabcolsep}{4pt}\footnotesize%
\begin{tabularx}{\textwidth}{@{}C{1.6cm} L{5.0cm} Y@{}}
\toprule
\headrow \textbf{Bab} & \textbf{Yang ditambahkan} & \textbf{Akibatnya} \\
\midrule
Bab~XVI & Pengujian prototipe dan presentasi &
Bunyi dan partikel yang telah dipasang diuji oleh pemain lain; ketidakcocokan
waktu dan kekuatan yang tidak terlihat penulisnya akan muncul di sana. \\
\bottomrule
\end{tabularx}}
\end{table}

Dengan bab ini, seluruh bagian yang terlihat dan terdengar dari level Anda
sudah lengkap. Yang tersisa bukan menambahkan bagian baru, melainkan menguji
apa yang sudah ada, dan itulah isi Bab~XVI.
